% =============================================================================
\section{Terrain Appearance}
\label{sec:texture}
% =============================================================================

\subsection{Why projection alone fails}

The obvious way to texture the ground is to project the panorama onto it. This
fails in two distinct ways.

First, \emph{grazing incidence}. Equirectangular rows compress with distance: at
\SI{80}{\meter} a single panorama row covers roughly \SI{12}{\meter} of ground
when the terrain is flat, against \SI{0.5}{\meter} when the horizon is anchored at
its true elevation. The result is radial smearing.

Second, and less obviously, \emph{upright content}. Projecting assumes the image
at elevation $\varphi$ is ground at radius $h/\tan(-\varphi)$. That holds for bare
terrain and fails completely for anything vertical: a \SI{0.5}{\meter} flower spans
a range of elevations, each landing at a different ground radius, so one flower
smears into a streak metres long. Rendered top-down this produces a radial
pinwheel. It is also double-counting, since the same flowers are placed as 3D
instances on top of the ground.

\subsection{Layered splat material}

Texturing is therefore a weighted blend of layers. A panorama layer supplies real
photographic colour where projection is trustworthy, gated by a visibility weight
that falls off near the nadir and the horizon and is withheld where a texel's own
panorama pixel is vegetation or built structure rather than ground. Synthetic
material layers, generated per region type, supply the rest.

\subsection{Real-material pattern texturing}

Where the panorama layer is withheld, we prefer real material over generated
material. Following pattern-based texturing~\cite{neyret:inria-00537511}, small
reference patches are extracted from the most solid interior of each region in the
panorama and flattened --- by fitting a local plane to the patch's own depth
points and reprojecting, which removes genuine oblique foreshortening rather than
only equirectangular stretch --- then assembled over the mesh's own triangles with
blended borders so the tiling does not read as faceted.

Reference patches are drawn from the object-removed panorama with real ground
colour restored from the original wherever foreground removal ate the ground
rather than an occluder. Where an occluder genuinely was removed, its footprint is
refilled by inpainting from the surrounding real ground, so what remains is
neither the occluder nor fabricated fill.
